\section{Local Stability, Basins, and Separatrices}

\begin{definition}[Stable state]
A state is \textbf{locally stable} if trajectories that start sufficiently
near it tend to remain near it or return toward it after small perturbations.
\end{definition}

\begin{definition}[Basin and separatrix]
The \textbf{basin of attraction} of a stable state is the set of initial
conditions that flow toward it. A \textbf{separatrix} is a boundary in state
space that divides trajectories heading toward different long-run outcomes.
\end{definition}

\begin{discussion}
These are not decorative definitions. They give a sharper interpretation of
behavioural failure. A person does not always fail because there is ``too much
temptation.'' Sometimes the system simply starts too deep inside the wrong
basin, or too far from the separatrix, for a small intervention to redirect it.
That is why timing and preparation matter so much: they determine whether the
system is being pushed while it is still near a boundary or only after it has
fallen deeply into a self-reinforcing region.
\end{discussion}

\begin{example}[Post-break study re-entry as a separatrix problem]
Suppose a student finishes lunch and returns to the desk. If the textbook is
open, the next problem is marked, and the phone is absent, the post-break state
may lie near the boundary between a study basin and a distraction basin. A
small cue can then redirect the trajectory. If the desk is cluttered and the
phone is already active, the same student may start well inside the distraction
basin. The difference is not moral worth. It is initial condition geometry.
\end{example}
